\documentclass[10pt,a4paper]{amsart}
\usepackage[margin=19mm]{geometry}
\usepackage{amsmath,amssymb,mathtools}
\usepackage[scr=rsfs,cal=euler]{mathalfa}
\usepackage{xfrac}
\usepackage[unicode,hidelinks,bookmarks=false]{hyperref}
\newcommand{\ie}{i.e.\ }
\newcommand{\cf}{cf.\ }
\newcommand{\resp}{respectively\ }
\newcommand{\Subsection}{\subsection}
\providecommand{\nobreakdash}{\nobreak\hbox{-}}
\setlength{\emergencystretch}{2em}
\multlinegap0pt
\usepackage{bm}
\usepackage{bbm}
\usepackage{dsfont}
\usepackage{xspace}
\usepackage{cancel}
\usepackage{mathtools}
\usepackage{amsthm}
\makeatletter
\@ifundefined{theorem}{\newtheorem{theorem}{Theorem}}{}
\@ifundefined{claim}{\newtheorem*{claim}{Claim}}{}
\@ifundefined{lemma}{\newtheorem{lemma}[theorem]{Lemma}}{}
\@ifundefined{proposition}{\newtheorem{proposition}[theorem]{Proposition}}{}
\@ifundefined{remark}{\newtheorem{remark}[theorem]{Remark}}{}
\makeatother
\newcommand{\hk}{\mathbf{HK}_{\Theta}}
\newcommand{\tot}{\overset{\Theta}{\to}}
\title[Homomorphisms of Hecke-Kiselman Monoids Associated to Simple]{Homomorphisms of Hecke-Kiselman Monoids Associated to Simple Oriented Graphs}
\author{Luka Andren\v{s}ek}
\date{Source: arXiv:2604.15497v3 (CC BY 4.0)}
\begin{document}
\maketitle

\noindent\emph{Source and licence.} Homomorphisms of Hecke-Kiselman Monoids Associated to Simple Oriented Graphs. Version arXiv:2604.15497v3 (CC BY 4.0), distributed under the Creative Commons Attribution 4.0 International licence, as recorded in the arXiv licence metadata for this exact version and shown on the abstract page https://arxiv.org/abs/2604.15497v3. Full rights notice: licenses/arxiv-2604.15497-CC-BY-4.0.txt.\par\smallskip

\noindent\textit{Editorial scope.} Counted unit: the Proposition whose source label is \texttt{prop2} in the article (source0.tex lines 902--905, source label \texttt{prop2}), delivered together with the article's own complete original proof of it (source0.tex lines 907--985, one \texttt{proof} environment of 79 lines). No mathematics is written, repaired or re-derived: statement and proof are the article's own text line for line. Required context retained uncounted: A-theta (lines 509--512); hk\_idempotents (lines 523--526); idempotent:formula (lines 545--552); IDS:no\_arrows (lines 632--638). Every definition, display and label of those lines is retained unchanged. Omitted from this package and not redistributed: 1--508, 513--522, 527--544, 553--631, 639--901, 906--906, 986--1534. Nothing inside a retained excerpt is omitted or altered: every retained body is delivered exactly as the article's lines are written, so no omission receipt of any kind accompanies this package. Citations: the retained excerpts carry no citation command at all, so no bibliography entry of the article is transcribed and no reference is redistributed. Reference adaptation: none was needed and none was made; every reference inside the delivered unit resolves inside it, so no unresolved reference or citation is produced. Printed slips kept literal: no printed word, symbol, hypothesis or number of the counted statement or of its proof was altered. Layout and class: the article's own class is \texttt{sn-jnl} with packages including fontenc, inputenc, bbm, amsmath, amsthm, amssymb, amsfonts, hyperref, enumitem, mathrsfs, tikz, graphicx; the wrapper is the lane's pinned \texttt{amsart} editorial wrapper at 10pt on A4, which restates the theorem-like environments the retained text needs and adds the article's own macro definitions that the retained text needs (\texttt{\textbackslash hk}, \texttt{\textbackslash tot}), with extra wrapper packages (bm, bbm, dsfont, xspace, cancel, mathtools). The delivered numbering is the unit's own. Assets: no retained span contains a figure environment, a graphics-inclusion command, a PDF-page inclusion, a table or a TikZ picture; no image file is copied into this package, the package declares no asset dependency and no image file is redistributed with it. Licence: version arXiv:2604.15497v3 of this article is distributed under the Creative Commons Attribution 4.0 International licence, as recorded in the arXiv licence metadata for this exact version and shown on the abstract page https://arxiv.org/abs/2604.15497v3. A full rights notice is delivered at licenses/arxiv-2604.15497-CC-BY-4.0.txt. Characters: every retained excerpt is pure ASCII, so the delivered engine, XeLaTeX with its default Latin Modern text font, reports no missing character.
\begin{remark}\label{A-theta}
   If $Y \subseteq X \in \mathcal{A}_{\Theta}$, then $\Theta[Y]$ is a subgraph of $\Theta[X]$, and hence it does not contain any oriented cycles.
   Therefore, $Y \in \mathcal{A}_{\Theta}$.
\end{remark}

\begin{proposition}\label{hk_idempotents}
    Let $X, Y \in \mathcal{A}_{\Theta}$.
    Then $e_X e_Y$ is an idempotent in $\hk$ if and only if $X \cup Y \in \mathcal{A}_{\Theta}$ and $e_X e_Y = e_{X \cup Y}$.
\end{proposition}

\begin{lemma}\label{idempotent:formula}
   Let $X \in \mathcal{A}_{\Theta}$.
   Then for any ordering of the set $X = \{x_{i_1}, x_{i_2}, \dots, x_{i_k}\}$ such that $x_{i_j} \tot x_{i_{j'}}$ implies $j < j'$ for all $j, j' \in \{1, 2, \dots, k\}$, we have
   \[
   e_X = x_{i_1} x_{i_2} \dots x_{i_k}.
   \]
   Moreover, we have $e_X z = z e_X = e_X$ for any $z \in \hk$ with $c(z) \subseteq X$.
\end{lemma}

\begin{lemma}\label{IDS:no_arrows}
   Let $X, Y \in \mathcal{A}_{\Theta}$.
   If there are no arrows starting in $Y$ and ending in $X$ and $X \cap Y = \emptyset$, then $X \cup Y \in \mathcal{A}_{\Theta}$ and
   \[
   e_X e_Y = e_{X \cup Y}.
   \]
\end{lemma}

\begin{proposition}\label{prop2}
   Let $X, Y \in \mathcal{A}_{\Theta}$.
   If $p_{\Theta}(X, Y)$ holds, then $e_X e_Y$ is an idempotent in $\hk$.
\end{proposition}

\begin{proof}
   Assume $p_{\Theta}(X, Y)$ holds.
   To prove that $e_X e_Y$ is an idempotent, by Proposition~\ref{hk_idempotents}, we need to show $X \cup Y \in \mathcal{A}_{\Theta}$ and $e_X e_Y = e_{X \cup Y}$.
   
   Denote $I = X \cap Y$ and
   \[
   R = \{x_i \in I \mid \exists x_b \in Y \setminus X, r \in \mathbb{N}_0,
   x_{i_1}, \dots, x_{i_r} \in I
   : x_b \tot x_{i_1} \tot \dots \tot x_{i_r} \tot x_i\}.
   \]
   Also set $S = I \setminus R$.
   We denote by
   \begin{equation}\label{M-N}
      M = (Y \setminus X) \cup R \quad \text{and} \quad N = (X \setminus Y) \cup S. 
   \end{equation}
   Observe that since $R, S \subseteq I = X \cap Y$, the above unions are disjoint.
   Then we further have 
   \begin{equation}\label{xy_decomp}
      X = N \cup R, \quad Y = M \cup S, \quad X \cup Y = N \cup M.
   \end{equation}
   The unions on the right-hand sides of these equalities are disjoint as well.
   The key observation is the following claim.
   \begin{claim}
      There are no arrows starting in $M$ and ending in $N$.
   \end{claim}

   \begin{proof}[Proof of the claim]
      We argue by contradiction.
      Let $x_i \in M$, $x_j \in N$ and suppose $x_i \tot x_j$.
      Now we consider cases on $x_i$ and $x_j$ based on~\eqref{M-N}.
      Since $p_{\Theta}(X, Y)$ holds, we cannot have $x_i \in Y \setminus X$ and $x_j \in X \setminus Y$.
      If $x_i \in Y \setminus X$ and $x_j \in S$, this implies $x_j \in R$ by the definition of $R$, which is a contradiction since $R \cap S = \emptyset$.
      If $x_i \in R$ and $x_j \in X \setminus Y$, then by the definition of $R$, there exist $x_b \in Y \setminus X$, $r \in \mathbb{N}_0$, and $x_{i_1}, \dots, x_{i_r} \in I$ such that
      \[
      x_b \tot x_{i_1} \tot \dots \tot x_{i_r} \tot x_i \tot x_j.
      \]
      Since $x_i \in R \subseteq I$ and $x_j \in X \setminus Y$, this contradicts the fact that $p_{\Theta}(X, Y)$ holds.
      If $x_i \in R$ and $x_j \in S$, we have some $x_b \in Y \setminus X$, $r \in \mathbb{N}_0$, and $x_{i_1}, \dots, x_{i_r} \in I$ such that 
      \[
      x_b \tot x_{i_1} \tot \dots \tot x_{i_r} \tot x_i \tot x_j.
      \]
      Since $x_i \in R \subseteq I$, this implies that $x_j \in R$.
      Thus $x_j \in R \cap S$, which is again a contradiction.
      The above cases are exhaustive by~\eqref{M-N}.
      Hence, in all cases, we have reached a contradiction, and therefore there are no arrows starting in $M$ and ending in $N$. 
   \end{proof}

   Since $N \subseteq X \in \mathcal{A}_{\Theta}$ and 
   $M \subseteq Y \in \mathcal{A}_{\Theta}$, Remark~\ref{A-theta} implies $N, M \in \mathcal{A}_{\Theta}$.
   Since $N \cap M = \emptyset$ and there are no arrows starting in $M$ and ending in $N$, Lemma~\ref{IDS:no_arrows} yields 
   \[
   N \cup M \in \mathcal{A}_{\Theta} \quad \text{and} \quad e_N e_M = e_{N \cup M}.
   \]
   Since $X \cup Y = N \cup M$ by \eqref{xy_decomp}, we obtain 
   \begin{equation}\label{t1}
      X \cup Y \in \mathcal{A}_{\Theta} \quad \text{and} \quad e_N e_M = e_{X \cup Y}.
   \end{equation}
   
   Since $R, S \subseteq X \cup Y \in \mathcal{A}_{\Theta}$, we also have $R, S \in \mathcal{A}_{\Theta}$ by Remark~\ref{A-theta}.
   As $R \subseteq M$ and there are no arrows starting in $M$ and ending in $N$, there are no arrows starting in $R$ and ending in $N$.
   Since $N \cap R = \emptyset$, Lemma~\ref{IDS:no_arrows} implies 
   \begin{equation}\label{t2}
      e_N e_R = e_{N \cup R} = e_X.
   \end{equation}
   Similarly, since $S \subseteq N$, there are no arrows starting in $M$ and ending in $S$.
   Since $S \cap M = \emptyset$, Lemma~\ref{IDS:no_arrows} implies 
   \begin{equation}\label{t3}
      e_S e_M = e_{S \cup M} = e_Y.
   \end{equation}
   Since $c(e_S) = S \subseteq X$ and $c(e_R) = R \subseteq M$, Lemma~\ref{idempotent:formula} implies 
   \begin{equation}\label{t4}
      e_X e_S = e_X \quad \text{and} \quad e_R e_M = e_M.
   \end{equation}
   Using \eqref{t1}, \eqref{t2}, \eqref{t3}, and \eqref{t4}, we obtain
   \[
   e_X e_Y = e_X e_S e_M = e_X e_M = e_N e_R e_M = e_N e_M = e_{X \cup Y}.
   \]
   This concludes the proof.
\end{proof}

\end{document}
